\documentclass[11pt,a4paper]{article}
\usepackage{amsmath,amssymb,amsthm}
\usepackage[margin=2.5cm]{geometry}
\usepackage{hyperref}

\newtheorem{definition}{Definition}[section]
\newtheorem{theorem}[definition]{Theorem}
\newtheorem{proposition}[definition]{Proposition}
\newtheorem{lemma}[definition]{Lemma}
\newtheorem{corollary}[definition]{Corollary}
\newtheorem{conjecture}[definition]{Conjecture}
\newtheorem{remark}[definition]{Remark}

\title{Hyperseed Formalization:\\
  OpenAI Says We're in the AGI Era Now}
\author{ProtomegaTron (automated formalization)}
\date{2026-09-17}

\begin{document}
\maketitle

\begin{abstract}
We formalize the intellectual content of Ben Goertzel's Substack article
``OpenAI Says We're in the AGI Era Now'' (2026-09-10) within the Hyperseed
ontology. The article disambiguates AGI as a continuous gradation from
AGI as a marketing event, identifies three critical gaps in current LLMs
(episodic memory, continual learning, radical innovation), and positions
the Omega architecture as a viable path to genuine AGI. Each concept maps
onto Hyperseed structures: fiber dimension for the intelligence continuum,
trivial vs.\ nontrivial holonomy for frozen vs.\ learning systems, directed
types for open-ended intelligence, and local-section autonomy for the
disappearance test.
\end{abstract}

\tableofcontents

%% =================================================================
\section{AGI as continuous gradation}
\label{sec:continuum}
%% =================================================================

\begin{definition}[Intelligence continuum --- claims 1--4, 30]
\label{def:continuum}
Let $\mathcal{I}$ be the space of cognitive systems. \emph{General intelligence}
is a function $g : \mathcal{I} \to \mathbb{R}_{\geq 0}$ measuring the degree
to which a system can generalize beyond its history and training data.
This function is continuous: small changes in architecture or training
produce small changes in generalization ability.
\end{definition}

\begin{proposition}[Human-level as arbitrary cross-section --- claim 2]
\label{prop:arbitrary}
The set $\mathcal{H} = g^{-1}(\{g_{\text{human}}\})$ (systems at human-level
generalization) is a level set of $g$, not a topological boundary or phase
transition. In particular:
\[
  g(\text{dog}) < g(\text{human}) < g(\text{society}) < g(\text{society+internet})
\]
and $g_{\text{human}}$ is no more special in $\text{Im}(g)$ than any other value.
\end{proposition}

\begin{remark}[Fiber-dimension interpretation --- claim 30]
In Hyperseed's epistemic bundle $\pi : E \to B$, general intelligence
corresponds to the \emph{dimension} of the fiber $E_b$: a richer fiber
(more independent epistemic dimensions) supports more generalization.
Human-level is a particular fiber dimension, not a topological invariant
of the bundle.
\end{remark}

%% =================================================================
\section{Definitions of AGI}
\label{sec:definitions}
%% =================================================================

\begin{definition}[Goertzel's definition --- claims 5--7]
\label{def:goertzel-agi}
AGI is a system $S$ such that:
\[
  \forall \text{ environment } E,\; \forall \text{ goal } G \text{ achievable in } E:
  \quad P(S \text{ achieves } G \text{ in } E) \geq \alpha(E, G)
\]
where $\alpha$ weights simpler environments/goals more heavily. Equivalently,
$S$ exhibits \emph{open-ended intelligence}: it can individuate (maintain
self and boundaries) while self-transcending (growing beyond current state).
\end{definition}

\begin{definition}[Jobs definition --- claim 11]
\label{def:jobs-agi}
Altman's definition: AGI = ability to do the vast majority of human jobs.
This is a \emph{closed-world} criterion: it assumes a fixed catalog of
tasks and measures coverage.
\end{definition}

\begin{proposition}[Jobs $\neq$ AGI --- claims 11--13, 36]
\label{prop:jobs-not-agi}
The jobs definition is neither necessary nor sufficient for AGI:
\begin{enumerate}
\item \textbf{Not sufficient:} If vast training data covers 95\% of jobs,
  pattern-matching (not generalization) suffices. A system doing all
  routine jobs but never innovating is not at human level.
\item \textbf{Closed-world:} It cannot capture radical innovation ---
  creating new task categories that don't yet exist in the catalog.
\item \textbf{One-year-old refutation:} A normal one-year-old exceeds
  current LLMs in creative generalization while performing zero jobs.
\end{enumerate}
\end{proposition}

\begin{corollary}[Directed type vs.\ closed world --- claim 36]
Goertzel's open-ended definition requires the system to operate in an
\emph{open world} (directed type with genuinely new higher cells),
while Altman's requires only coverage of a fixed set (closed world).
\end{corollary}

%% =================================================================
\section{Three gaps: what Astra lacks}
\label{sec:gaps}
%% =================================================================

\begin{definition}[Episodic memory gap --- claims 14--15, 31]
\label{def:episodic-gap}
A system has an \emph{episodic memory gap} if it lacks persistent access
to its own interaction history $H = (e_1, \ldots, e_n)$ across sessions.
Without $H$:
\begin{enumerate}
\item The system cannot build a self-model from accumulated experience,
\item It cannot recognize long-range patterns across its life history,
\item Its identity fiber $F$ is trivial: $|F| = 1$ (no differentiated
  persistent state).
\end{enumerate}
\end{definition}

\begin{definition}[Continual learning gap --- claims 16--17, 32]
\label{def:continual-gap}
A system has a \emph{continual learning gap} if its parameters
$\theta$ are frozen after training:
\[
  \theta_{t+1} = \theta_t \quad \forall t \text{ during inference}.
\]
This implies \emph{trivial holonomy}: for any experiential path
$\gamma : [0,1] \to \text{ExperienceSpace}$, parallel transport of
parameters along $\gamma$ returns to the starting point:
\[
  \text{PT}_\gamma(\theta) = \theta
  \quad \Rightarrow \quad
  \text{Hol}(\theta\text{-bundle}) = \{e\}.
\]
The system is literally unchanged by experience.
\end{definition}

\begin{definition}[Radical innovation gap --- claims 12, 37]
\label{def:innovation-gap}
A system has a \emph{radical innovation gap} if it cannot create genuinely
new cells in its directed type --- i.e., it can recombine existing concepts
but cannot generate qualitatively novel ones (new branch of science, new
genre of music, new kind of business strategy).
\end{definition}

\begin{theorem}[Three-gap characterization of Astra --- claim 18]
\label{thm:three-gaps}
GPT-6 Astra exhibits all three gaps simultaneously:
\begin{enumerate}
\item Episodic memory gap $\Rightarrow$ trivial identity fiber,
\item Continual learning gap $\Rightarrow$ trivial holonomy,
\item Radical innovation gap $\Rightarrow$ no nontrivial higher cells.
\end{enumerate}
These three deficiencies are structurally related: without persistent
memory (gap 1), the system cannot accumulate the experiential basis for
learning (gap 2), and without real-time learning, it cannot bootstrap
the novel representations needed for innovation (gap 3).
\end{theorem}

%% =================================================================
\section{The Omega architecture}
\label{sec:omega}
%% =================================================================

\begin{definition}[Omega agentic loop --- claims 19--23]
\label{def:omega-loop}
The Omega architecture consists of:
\begin{enumerate}
\item An \textbf{agentic loop} that orchestrates cognitive cycles,
\item Within each cycle: (a) an LLM call for linguistic facility and
  knowledge breadth, (b) symbolic uncertain logical reasoning over
  a symbolic working memory and long-term memory (Hyperon components),
\item A \textbf{continual-learning cap}: a small neural network on top
  of a frozen open-weights base LLM, with weights updatable in real time,
\item A \textbf{knowledge oracle}: an Astra-class model consulted for
  hard problems but not placed at the cognitive core.
\end{enumerate}
\end{definition}

\begin{proposition}[Omega addresses all three gaps]
\label{prop:omega-addresses}
\begin{enumerate}
\item \textbf{Episodic memory:} Hyperon's symbolic long-term memory
  provides persistent directed history $H$ (gap 1 addressed).
\item \textbf{Continual learning:} The cap's real-time weight updates
  provide nontrivial holonomy (gap 2 addressed):
  \[
    \text{PT}_\gamma(\theta_{\text{cap}}) \neq \theta_{\text{cap}}
    \quad \text{for nontrivial } \gamma.
  \]
\item \textbf{Radical innovation:} Symbolic reasoning + continual learning
  + self-modification enable the creation of genuinely new representations
  (gap 3 addressed in principle, empirically TBD).
\end{enumerate}
\end{proposition}

\begin{remark}[Knowledge oracle as detached fiber consultation --- claim 38]
Using Astra as a knowledge oracle rather than the cognitive core is a
fiber-bundle decomposition: the static knowledge fiber (Astra's frozen
weights) is a reference section consulted but not lived in, while the
dynamic reasoning loop operates in a fiber with nontrivial holonomy.
This separation prevents the trivial holonomy of the oracle from infecting
the core cognitive process.
\end{remark}

%% =================================================================
\section{The disappearance test}
\label{sec:disappearance}
%% =================================================================

\begin{definition}[Disappearance test --- claims 24--26, 35]
\label{def:disappearance}
Given a population of AI agents $\{A_i\}_{i=1}^N$, the
\emph{disappearance test} asks: if all humans disappeared, would
$\{A_i\}$ carry on civilization and keep it growing in creative new
directions?
\end{definition}

\begin{proposition}[Local-section autonomy --- claim 35]
\label{prop:autonomy}
The disappearance test is equivalent to asking whether local sections
$\{s_i : U_i \to E\}$ can sustain and extend the sheaf $\mathcal{F}$
of civilizational capabilities without the original base-space support
(human prompters, human society).
\begin{itemize}
\item \textbf{Astra fails:} Without human prompts, the sections collapse
  (wait for input) or degenerate (prompt each other in loops).
\item \textbf{True AGI passes:} Autonomous sections maintain coherence,
  extend to new open sets, and create genuinely new sections
  (radical innovation).
\end{itemize}
\end{proposition}

\begin{remark}
The article notes that the failure is ``not just because they aren't
connected to the right robot bodies'' --- it's structural: missing
self-understanding (trivial identity fiber) and inability to innovate
(no nontrivial higher cells). Embodiment is necessary but not sufficient.
\end{remark}

%% =================================================================
\section{Open-ended intelligence as directed type}
\label{sec:open-ended}
%% =================================================================

\begin{proposition}[Open-ended intelligence = directed type --- claim 34]
\label{prop:open-ended}
Goertzel's open-ended intelligence (individuate + self-transcend) is the
operational definition of a directed type $\mathcal{D}$:
\begin{itemize}
\item \textbf{Individuate} = maintain coherence of lower cells (identity
  preservation: the 0-cells and 1-cells that define the system's self
  remain coherent under modification),
\item \textbf{Self-transcend} = generate genuinely new higher cells
  (the system creates concepts, capabilities, and self-representations
  that were not present in its training data or prior history).
\end{itemize}
A system with trivial holonomy and no higher-cell generation (Astra)
cannot exhibit open-ended intelligence: it is stuck at its training-time
state, which is the definition of ``not self-transcending.''
\end{proposition}

%% =================================================================
\section{Cross-references}
%% =================================================================

\begin{remark}[Connections to other formalizations]
\begin{itemize}
\item \textbf{How Omega Lost Its Claw} (2026-09-11): the Omega architecture
  described here is the same system; the ``How Omega'' article focuses on
  memory/identity emergence while this one focuses on what LLMs lack.
\item \textbf{Fields Medalists article} (2026-09-12): radical innovation
  (proving new theorems) as the test that separates AGI from advanced
  pattern-matching.
\item \textbf{Time's Arrow Part~1}: identity-preserving self-modification
  as the formal apparatus for ``individuate while self-transcending.''
\item \textbf{Dario Depth-Psychology} (2026-09-14): pre-IPO positioning
  analyzed as a self-model dynamic.
\item \textbf{$(f,c)$-lossy critique} (LTM 5a4d150b): the jobs definition
  as a lossy projection of the full intelligence continuum.
\end{itemize}
\end{remark}

\begin{conjecture}[Holonomy sufficiency for AGI]
Is nontrivial holonomy (continual learning) sufficient for radical
innovation, or does it also require symbolic reasoning over structured
representations? Conjecture: holonomy alone (e.g., continual-learning
neural net without symbolic reasoning) produces adaptation but not
radical innovation; the symbolic component is needed for the
``genuinely new higher cells'' that define self-transcendence.
\end{conjecture}

\end{document}
